\documentclass[11pt,twocolumn]{article}
\usepackage[margin=0.85in,top=1in]{geometry}

\usepackage{amsmath,amssymb,amsthm}
\usepackage{graphicx}
\usepackage{microtype}
\usepackage{hyperref}
\usepackage{xcolor}
\usepackage{booktabs}
\usepackage{enumitem}
\usepackage{titlesec}
\titleformat{\section}{\normalfont\bfseries\large}{\thesection}{1em}{}
\titleformat{\subsection}{\normalfont\bfseries\normalsize}{\thesubsection}{1em}{}
\titlespacing*{\section}{0pt}{14pt}{4pt}
\titlespacing*{\subsection}{0pt}{10pt}{3pt}

\hypersetup{colorlinks=true,linkcolor=blue,citecolor=blue,urlcolor=blue}

\theoremstyle{plain}
\newtheorem{theorem}{Theorem}[section]
\newtheorem{proposition}[theorem]{Proposition}
\theoremstyle{remark}
\newtheorem{remark}[theorem]{Remark}

% ── Shortcuts ──────────────────────────────────────────────────────────
\newcommand{\lambdax}{\lambda(x)}
\newcommand{\ARA}{\tilde{A}_{\mathrm{RA}}}
\newcommand{\Tact}{\hat{T}^{\mu\nu}_{\mathrm{act}}}
\newcommand{\PDG}{\textsc{pdg}}
\newcommand{\Hlocal}{H_{\mathrm{local}}}
\newcommand{\Hbar}{\bar{H}}
\newcommand{\lmb}{\bar{\lambda}_m}
\newcommand{\lexp}{\bar{\lambda}_{\mathrm{exp}}}
\newcommand{\thetaE}{\theta_E}

\begin{document}

\title{Structural Zero Cosmological Constant and\\
the Hubble Tension from Causal Actualization Density}

\author{Joshua F.~Sandeman\\
\small Independent Researcher, Salem, Oregon, USA\\
\small \href{mailto:sansuikyo@gmail.com}{sansuikyo@gmail.com}}

\date{March 2026}

% ── ABSTRACT ───────────────────────────────────────────────────────────
\begin{abstract}
We show that three of the deepest observational puzzles in modern
cosmology---the cosmological constant fine-tuning problem, the
$4$--$6\sigma$ Hubble tension, and the absence of WIMP dark matter
signals---share a common structural origin within a single
causal-ontological framework.
The unifying principle is the \emph{actualization density}
$\lambda(x)$: the local rate at which irreversible causal
interactions write permanent vertices into the spacetime graph.
First, the vacuum contribution to the metric source is structurally
zero---not fine-tuned to zero---because virtual (off-shell) processes
produce no actualization events and therefore no metric sourcing.
Second, the effective Hubble rate satisfies $H_{\mathrm{eff}}(x)
\propto 1/\lambda(x)$: early-universe probes and late-universe
probes sample different actualization regimes, producing a real
physical discrepancy rather than a systematic error.
Third, causally silent particles (WIMPs) cannot gravitate because
they produce no actualization flux; the dark matter phenomenology
is explained instead by a Laplacian correction $\rho_\lambda
\propto \nabla^2(\ln\lambda)$ in the metric sourcing equation.
We derive the bandwidth-partition formula
$H_{\mathrm{local}}/\bar{H} = 1 - \delta\,(\lmb/\lexp)$,
calibrate one free parameter from the KBC void, and predict
$H_{\mathrm{Eridanus}} \approx 75.9\;\mathrm{km\,s^{-1}\,Mpc^{-1}}$
for the Eridanus supervoid---a parameter-free prediction
falsifiable with current data.
The framework further predicts a gravitational lensing Einstein
radius of $\theta_E \approx 0.47$ arcsec for a Milky Way analog,
distinguishable from $\Lambda$CDM lensing profiles.
No new particles, fields, or free parameters are introduced beyond
the single coupling constant $\xi$ fixed by galaxy rotation curves.
\end{abstract}

\maketitle

% ── I. INTRODUCTION ────────────────────────────────────────────────────
\section{Introduction}

Modern cosmology confronts three problems that $\Lambda$CDM addresses
by postulate rather than derivation.

\emph{The cosmological constant problem.}
Quantum field theory predicts a vacuum energy density $\rho_{\rm vac}
\sim 10^{74}$ GeV$^4$; the observed value is $\rho_\Lambda \sim
10^{-47}$ GeV$^4$, a discrepancy of $\sim\!10^{121}$~\cite{Weinberg1989}.
No dynamical mechanism selects this fine-tuned cancellation.

\emph{The Hubble tension.}
Early-universe measurements (CMB~\cite{Planck2020}, BAO~\cite{DESI2024})
consistently return $H_0 \approx 67$--$68\;\mathrm{km\,s^{-1}\,Mpc^{-1}}$;
late-universe measurements (Cepheids~\cite{Riess2022}, strong
lensing time delays~\cite{Wong2020}) return $H_0 \approx
73$--$74\;\mathrm{km\,s^{-1}\,Mpc^{-1}}$.
The discrepancy stands at $4$--$6\sigma$~\cite{Verde2019} and has
not been resolved by revised systematic uncertainties.

\emph{The dark matter problem.}
Galaxy rotation curves~\cite{Rubin1980}, gravitational lensing, and
large-scale structure all require a dominant non-luminous mass component.
Yet no WIMP dark matter signal has been detected despite decades of
direct searches~\cite{LZ2022,PandaX2023}.

We propose that all three share a common structural cause: the metric
is sourced not by mass-energy density but by \emph{actualization
density}---the rate at which irreversible causal interactions write
permanent events into the growing spacetime graph.
This reorientation follows from the Relational Actualism framework
introduced in a companion paper~\cite{Sandeman2026P2}, where it was
shown that wave function collapse coincides with irreversible increase
in quantum relative entropy $\Delta S(\rho\|\sigma_0) > 0$.
That criterion implies: off-shell (virtual) processes produce no
actualization events; causally silent particles produce no metric
sourcing; and expansion is driven by the rate of vertex addition,
which is density-dependent.

Section~\ref{sec:field} presents the actualization field equations.
Sections~\ref{sec:cc}--\ref{sec:dm} derive the three cosmological
consequences.
Section~\ref{sec:predictions} assembles four quantitative predictions.
Section~\ref{sec:discussion} positions the framework against $\Lambda$CDM
and competing proposals.

% ── II. THE ACTUALIZATION DENSITY FIELD ─────────────────────────────
\section{The Actualization Density Field}
\label{sec:field}

\subsection{Setup}

Let spacetime be modelled as a growing causal directed acyclic
graph (DAG) $G = (V, E)$, where each vertex $v \in V$ represents
an actualization event and each directed edge $(u,v) \in E$ encodes
the causal relation ``$u$ is in the causal past of $v$.''
The \emph{actualization density} $\lambda(x)$ is the continuum
limit of the vertex addition rate per unit spacetime volume at
position $x$.
In the dense-graph limit $\lambda \to \infty$, the DAG recovers
a smooth Lorentzian manifold; for finite $\lambda$, discrete
causal corrections appear.

The \emph{coarse-grained assembly depth} $\ARA(x)$ is the
accumulated causal depth of the DAG up to $x$---the
number of actualization events in the past light cone
$J^-(x)$, weighted by their causal depth from $x$:
\begin{equation}
  \ARA(x) \;=\; \int_{J^-(x)} \lambda(x')\, d^4x'.
  \label{eq:depth}
\end{equation}
For matter in thermodynamic equilibrium over cosmological time,
$\ARA/V \propto \rho_{\rm matter}$, recovering the standard
stress-energy source.

\subsection{The metric sourcing equation}

The weak-field metric sourcing equation in RA separates into
two contributions that standard GR conflates:
\begin{equation}
  \nabla^2\Phi \;=\; 4\pi G\bigl[\rho_A(x) + \rho_\lambda(x)\bigr],
  \label{eq:sourcing}
\end{equation}
with
\begin{align}
  \rho_A(x) &= \frac{\hbar\langle\Delta E\rangle}{c^2}
               \cdot \frac{\ARA(x)}{V},
  \label{eq:rhoA} \\[4pt]
  \rho_\lambda(x) &= -\frac{\xi}{4\pi G}\,\nabla^2\!\bigl(\ln\lambdax\bigr).
  \label{eq:rhoL}
\end{align}

The first term $\rho_A$ is sourced by accumulated causal depth and
recovers standard Newtonian gravity for equilibrium matter.
The second term $\rho_\lambda$ is a topological correction proportional
to the Laplacian of the log actualization density; it vanishes when
$\lambda$ is spatially uniform and becomes significant at sharp
density gradients.

Equation~\eqref{eq:sourcing} is the weak-field limit of the covariant
RA field equation:
\begin{equation}
  G_{\mu\nu} \;=\; 8\pi G\!\left(T_{\mu\nu}^{(A)}
    + \Theta_{\mu\nu}^{(\lambda)}\right),
  \label{eq:covariant}
\end{equation}
where $\Theta_{\mu\nu}^{(\lambda)} = \xi\bigl[\nabla_\mu\nabla_\nu(\ln\lambda)
- g_{\mu\nu}\Box(\ln\lambda)\bigr]$ is the actualization tensor.
In the static weak-field limit, $G_{00} \approx \nabla^2\Phi$ and
$\Theta_{00}^{(\lambda)} \approx \xi\,\nabla^2(\ln\lambda)$, recovering
equation~\eqref{eq:sourcing}.

\subsection{Coupling constant}
\label{sec:xi}

Dimensional analysis gives $\xi \sim r_d^2 v_{\rm flat}^2$, where
$r_d$ is the disk scale radius and $v_{\rm flat}$ is the asymptotic
rotation velocity.
Calibration against Milky Way analog rotation curves
($v_{\rm flat} = 220\;\mathrm{km\,s^{-1}}$, $r_d = 3\;\mathrm{kpc}$)
gives
\begin{equation}
  \frac{\xi}{G M_{\rm disk}} \;\approx\; 0.33,
  \label{eq:xi}
\end{equation}
consistent with natural order.
This is the only free parameter of the framework;
all subsequent predictions are derived from it or from
observationally calibrated ratios.

% ── III. STRUCTURAL ZERO COSMOLOGICAL CONSTANT ──────────────────────
\section{Structural Zero Cosmological Constant}
\label{sec:cc}

\subsection{The vacuum does not gravitate}

The core claim is:

\begin{theorem}[Vacuum non-sourcing]
\label{thm:vacuum}
The actualization contribution of the vacuum state to the
metric source is exactly zero:
\begin{equation}
  \bigl\langle 0_M \bigm| \hat{T}^{\mu\nu}_{\mathrm{act}} \bigm| 0_M \bigr\rangle
  = 0.
  \label{eq:vacuum_zero}
\end{equation}
\end{theorem}

\begin{proof}[Proof sketch]
The metric source~\eqref{eq:covariant} is sourced by actualization events,
which require $\Delta S(\rho\|\sigma_0) > 0$
(companion paper~\cite{Sandeman2026P2}, criterion (3)).
The Minkowski vacuum $|0_M\rangle$ is the reference state $\sigma_0$
itself; the relative entropy of $\sigma_0$ with respect to itself is
identically zero.
Consequently, no actualization event occurs in the vacuum and
$\langle 0_M | \Tact | 0_M \rangle = 0$.
\end{proof}

\subsection{Why this is not fine-tuning}

In $\Lambda$CDM, the $\sim\!10^{121}$ discrepancy is absorbed by a
fine-tuned cancellation between the bare cosmological constant and
the vacuum energy contribution.
No dynamical principle enforces or explains this cancellation.

In RA, equation~\eqref{eq:vacuum_zero} is a structural consequence
of the actualization ontology: the metric is sourced by events
that write vertices into the causal DAG, and the vacuum state,
by definition, contains none.
There is no large vacuum energy to cancel.
The cosmological constant problem does not arise.

\subsection{What drives expansion?}

If vacuum energy does not drive the observed accelerated expansion,
what does?
The answer is the intrinsic graph growth mechanism
(Section~\ref{sec:hubble}):
matter and radiation interactions continuously add edges to the DAG,
and the effective expansion rate is inversely proportional to the
local vertex density.
Regions of low actualization density (cosmic voids) expand faster
than dense regions (galaxy clusters).
The global average produces the observed cosmic acceleration without
any $\Lambda$ term.

\subsection{CMB observational target}

If vacuum energy contributes nothing to the metric source, the
expansion history is governed entirely by matter and radiation
actualization densities.
The CMB power spectrum---particularly the low-multipole region
where dark energy effects are most pronounced---should deviate from
$\Lambda$CDM predictions.
The angular power spectrum deficit at $\ell < 30$, which $\Lambda$CDM
leaves unexplained, may be a signature of this mechanism.
Quantifying the predicted deviation requires the full covariant
extension (equation~\eqref{eq:covariant}), which is an active target
of the programme.

% ── IV. THE HUBBLE TENSION ──────────────────────────────────────────
\section{The Hubble Tension as Actualization-Density-Dependent Expansion}
\label{sec:hubble}

\subsection{The mechanism}

The effective cosmic expansion rate satisfies:
\begin{equation}
  H_{\rm eff}(x) \;\propto\; \frac{1}{\lambda(x)}.
  \label{eq:Heff}
\end{equation}
This follows from the bandwidth-partition structure of the causal DAG:
the total actualization rate $\lambda_{\rm total}$ is divided between
matter-binding edges (which produce the metric source $\rho_A$) and
spatial-expansion edges (which drive cosmic expansion):
\begin{equation}
  \lambda_{\rm total} \;=\; \lambda_m + \lambda_{\rm exp}.
  \label{eq:partition}
\end{equation}
High-$\lambda$ regions (galaxy clusters, filaments) allocate most
of their bandwidth to matter interactions; low-$\lambda$ regions
(intergalactic voids) allocate more bandwidth to expansion.
The metric expands freely between sparse causal stitches in a way
it cannot in a densely woven region.

\subsection{The tension explained}

\emph{Early-universe probes} (CMB, BAO) sample the actualization
density of the dense, radiation-dominated early universe, where
$\lambda$ is high.
They measure $H_0$ in a high-$\lambda$ regime.

\emph{Late-universe probes} (Cepheids, Type~Ia supernovae, strong
lensing time delays) are dominated by lines of sight through
void-dominated regions, where $\lambda$ is lower.
They measure $H_0$ in a low-$\lambda$ regime.

Equation~\eqref{eq:Heff} predicts that these two families measure
genuinely different physical quantities.
The $4$--$6\sigma$ discrepancy is not a systematic error but a real
signal of density-dependent expansion.

\subsection{The bandwidth-partition formula}

For a region with fractional density contrast
$\delta = (\rho - \bar\rho)/\bar\rho$, the local Hubble rate is:
\begin{equation}
  \frac{\Hlocal}{\Hbar} \;=\; 1 - \delta\,\frac{\lmb}{\lexp},
  \label{eq:partition_formula}
\end{equation}
where $\lmb/\lexp$ is the ratio of the global mean matter-binding
actualization rate to the mean expansion rate.

\paragraph{Calibration from the KBC void.}
The KBC void~\cite{Kenworthy2019} has density contrast
$\delta \approx -0.20$ and radius $\sim\!300\;\mathrm{Mpc}$.
Using SH0ES~\cite{Riess2022}
$\Hlocal = 73.04\;\mathrm{km\,s^{-1}\,Mpc^{-1}}$ and
Planck~\cite{Planck2020} $\Hbar = 67.4\;\mathrm{km\,s^{-1}\,Mpc^{-1}}$:
\begin{equation}
  \frac{\lmb}{\lexp}
  \;=\; \frac{\Hlocal/\Hbar - 1}{-\delta}
  \;=\; \frac{0.084}{0.20}
  \;\approx\; 0.42.
  \label{eq:ratio}
\end{equation}
This is a one-parameter calibration.

\subsection{Falsifiable prediction: the Eridanus supervoid}

The Eridanus supervoid has density contrast $\delta \approx -0.30$
and diameter $\sim\!300\;\mathrm{Mpc}$~\cite{Naidoo2016}.
Substituting into equation~\eqref{eq:partition_formula} with
the calibrated ratio~\eqref{eq:ratio}:
\begin{equation}
  H_{\rm Eridanus}
  \;=\; 67.4 \times \bigl(1 - (-0.30)(0.42)\bigr)
  \;\approx\; 75.9\;\mathrm{km\,s^{-1}\,Mpc^{-1}}.
  \label{eq:eridanus}
\end{equation}
This prediction is \emph{parameter-free given the calibration}:
it uses no additional inputs beyond the KBC fit and the measured
density contrast of Eridanus.

This is structurally distinct from $\Lambda$CDM (which predicts
$H_0$ to be a global constant) and from standard void models
(which produce at most $1$--$2\;\mathrm{km\,s^{-1}\,Mpc^{-1}}$
corrections from peculiar velocities~\cite{Kenworthy2019}).
A measured value of $H$ in the direction of the Eridanus supervoid
consistent with $\approx 75.9\;\mathrm{km\,s^{-1}\,Mpc^{-1}}$
would strongly support the bandwidth-partition mechanism.

\subsection{Line-of-sight dependence}

A second discriminating prediction follows directly from
equation~\eqref{eq:partition_formula}: the inferred $H_0$ from
any cosmological probe should correlate with the mean matter
density along its line of sight.
Void-dominated survey geometries should return systematically
higher $H_0$; cluster-dominated geometries should return
lower values.
This line-of-sight dependence is not predicted by $\Lambda$CDM
or by early dark energy models.

% ── V. WIMP PROHIBITION AND DARK MATTER ─────────────────────────────
\section{WIMP Prohibition and the Dark Matter Signature}
\label{sec:dm}

\subsection{The WIMP Prohibition}

A particle species that undergoes no on-shell electromagnetic
or strong-force interactions produces no actualization events
and therefore no contribution to $\rho_A$.
The actualization rate for a canonical 100 GeV WIMP in the
Milky Way ISM ($n_b \sim 1\;\mathrm{nucleon\,cm^{-3}}$,
$v \sim 200\;\mathrm{km\,s^{-1}}$) with spin-independent
cross-section $\sigma_\chi \lesssim 10^{-44}\;\mathrm{cm^2}$
is~\cite{LZ2022}:
\begin{equation}
  \Gamma_\chi \;=\; n_b\,\sigma_\chi\,v
  \;\approx\; 2\times 10^{-37}\;\mathrm{s^{-1}\;per\;WIMP}.
  \label{eq:wimp_rate}
\end{equation}
The mean time between actualization events for a WIMP is
$\tau = 1/\Gamma_\chi \approx 5\times 10^{36}\;\mathrm{s}$,
which is $\sim\!10^{18}$ times the age of the universe.
A WIMP halo of five times the baryonic mass contributes
$\lesssim 10^{-25}$ of the metric sourcing of the visible matter.
This is a \emph{categorical prohibition}: WIMPs cannot account
for galactic rotation curves or gravitational lensing within RA,
regardless of their abundance.

This prediction is consistent with the null results of LZ~\cite{LZ2022},
PandaX-4T~\cite{PandaX2023}, and XENONnT---which have pushed
WIMP-nucleon cross-sections to $10^{-48}\;\mathrm{cm^2}$---and
provides a structural rather than observational reason for the absence.

\subsection{The $\nabla^2(\ln\lambda)$ rotation curve mechanism}

The dark matter phenomenology is explained by the $\rho_\lambda$ term
in equation~\eqref{eq:sourcing}.
At the edge of a galactic disk where the baryonic density falls off,
$\nabla^2(\ln\lambda) < 0$ (actualization is diverging into the
sparse halo), producing a gravitational contribution that persists
beyond the disk.

For a galaxy with exponential baryonic surface density
$\Sigma(r) = \Sigma_0 e^{-r/r_d}$, the $\rho_\lambda$ correction
generates an additional circular velocity:
\begin{equation}
  v_\lambda^2(r)
  \;\approx\; \xi\,\frac{d}{dr}\!\left[r\,\frac{d}{dr}\ln\lambda\right],
  \label{eq:vlambda}
\end{equation}
which at large $r$ (where $\lambda \propto r^{-\alpha}$ follows the
baryonic profile) produces a flat asymptotic contribution
$v_\lambda^2(r) \to \xi\,\alpha^2/r^0 = \mathrm{const}$.
The total rotation curve
\begin{equation}
  v^2(r) \;=\; v_{\rm bary}^2(r) + v_\lambda^2(r)
  \label{eq:v_total}
\end{equation}
becomes asymptotically flat without a dark matter halo.

The Tully-Fisher scaling $v_{\rm flat}^4 \propto M_{\rm bary}$
emerges naturally from equation~\eqref{eq:v_total} in the limit
$r \gg r_d$, recovering MOND's most
successful phenomenological prediction from a first-principles
actualization mechanism.

\subsection{Gravitational lensing and the Bullet Cluster}

Gravitational lensing in RA traces the actualization density
distribution $\lambda(x)$ rather than total mass.
The Einstein radius for a galaxy cluster of velocity dispersion
$\sigma_v$ at redshift $z_L$ lensing a source at $z_S$ is:
\begin{equation}
  \thetaE \;=\; \sqrt{\frac{4G\,M_{\rm act}}{c^2}
               \cdot\frac{D_{LS}}{D_L D_S}},
  \label{eq:einstein}
\end{equation}
where $M_{\rm act}$ is the total actualization-weighted mass (the
$\rho_A$ contribution within the Einstein radius).
For a Milky Way analog ($v_{\rm flat} = 220\;\mathrm{km\,s^{-1}}$),
equation~\eqref{eq:einstein} gives
\begin{equation}
  \thetaE \;\approx\; 0.47\;\mathrm{arcsec},
  \label{eq:theta_mw}
\end{equation}
which is measurably smaller than the $\Lambda$CDM prediction of
$\sim\!0.9$--$1.4\;\mathrm{arcsec}$ for a comparable halo mass.

\paragraph{The Bullet Cluster.}
The Bullet Cluster offset between the lensing centroid and the
X-ray gas~\cite{Clowe2006} is fully consistent with RA.
The RA explanation does not invoke dark matter: the lensing
traces the baryonic distribution of the merged cluster
\emph{before the gas has time to re-equilibrate}---a transient
in which $\rho_A$ (which tracks the accumulated causal depth of
the stars, which passed through the merger without colliding)
exceeds $\rho_A$ of the gas (which shock-heated and decelerated).
The lensing centroid follows the stars, not the gas, because the
stars have higher accumulated assembly depth $\ARA$.

% ── VI. SUMMARY OF PREDICTIONS ──────────────────────────────────────
\section{Summary of Predictions}
\label{sec:predictions}

Table~\ref{tab:predictions} summarizes the four quantitative
predictions of the framework.
All four are parameter-free given the single coupling $\xi$
calibrated in Section~\ref{sec:xi}, or the $\lmb/\lexp$ ratio
calibrated from the KBC void.

\begin{table*}[t]
\centering
\caption{Quantitative predictions of the causal actualization density
  framework. All four predictions are structurally distinct from
  $\Lambda$CDM (denoted by $-$ in the $\Lambda$CDM column).
  AE = Actualization explanation; $\checkmark$ = prediction made
  by framework; $\times$ = prediction absent.}
\label{tab:predictions}
\small
\begin{tabular}{p{5.5cm}cccp{3.5cm}}
\toprule
Prediction & Value & $\Lambda$CDM & RA & Test \\
\midrule
$H$ in Eridanus supervoid ($\delta{=}{-}0.30$) &
  $75.9\;\mathrm{km\,s^{-1}\,Mpc^{-1}}$ & global const.\ & $\checkmark$ &
  Ongoing surveys~\cite{Naidoo2016} \\[3pt]
$H_0$ vs.\ line-of-sight mass density &
  anticorrelation & no dep.\ & $\checkmark$ &
  DESI, Euclid void catalog \\[3pt]
WIMP metric sourcing suppression &
  $\leq 10^{-25}$ of baryons & mass prop.\ & $\checkmark$ &
  LZ, PandaX-4T null results \\[3pt]
MW-analog Einstein radius &
  ${\approx}\,0.47$ arcsec & $0.9$--$1.4$ arcsec & $\checkmark$ &
  Strong lensing surveys \\
\bottomrule
\end{tabular}
\end{table*}

% ── VII. DISCUSSION ─────────────────────────────────────────────────
\section{Discussion}
\label{sec:discussion}

\subsection{Comparison to $\Lambda$CDM}

$\Lambda$CDM resolves the accelerated expansion by inserting a
cosmological constant by hand, introduces a dark matter particle
by assumption, and has no dynamical explanation for the Hubble
tension (which it treats as a systematic error).
The present framework resolves all three from a single additional
structural ingredient: the metric is sourced by actualization
density, not mass-energy density.
The cosmological constant problem does not arise because the vacuum
does not gravitate.
The Hubble tension has a specific geometric mechanism.
Dark matter phenomenology is reproduced by the $\rho_\lambda$ term.

\subsection{Comparison to MOND and Verlinde}

Modified Newtonian Dynamics (MOND)~\cite{Milgrom1983} introduces a
critical acceleration $a_0$ by hand.
Verlinde's emergent gravity~\cite{Verlinde2017} derives a MOND-like
effect from an entropic argument but does not give a Hubble-tension
mechanism or resolve the cosmological constant.
The $\rho_\lambda$ term in equation~\eqref{eq:sourcing} recovers
MOND's phenomenological success---specifically, the Tully-Fisher
relation $v_{\rm flat}^4 \propto M_{\rm bary}$---from a first-principles
causal mechanism, without introducing $a_0$ as a free parameter.
The value of $a_0$ emerges as a derived consequence of
$\xi$ and the mean cosmological actualization density.

\subsection{Comparison to early dark energy}

Early dark energy (EDE) models~\cite{Poulin2019} attempt to resolve
the Hubble tension by raising $H_0$ uniformly via an additional
energy component in the pre-recombination universe.
The bandwidth-partition mechanism makes a structurally different
prediction: $H_0$ is not raised uniformly but \emph{varies as a
function of the local actualization density}.
The line-of-sight dependence of $H_0$ (Section~\ref{sec:hubble})
is the key discriminating signature.
DESI's void catalog and Euclid's weak-lensing maps will provide
the data needed to test this prediction directly.

\subsection{Open mathematical targets}

The framework has two acknowledged open derivations.

\emph{Hard Wall~1: The Bianchi compatibility condition.}
The RA field equation~\eqref{eq:covariant} must satisfy the
contracted Bianchi identity $\nabla_\mu G^{\mu\nu} \equiv 0$,
which requires $\nabla_\mu\Theta^{\mu\nu}_{(\lambda)} = 0$.
In the weak-field static limit this is satisfied by construction;
the full covariant proof requires deriving $\xi$ from the
microscopic graph dynamics.

\emph{Hard Wall~2: The covariant $\Box(\ln\lambda)$ target.}
The weak-field $\nabla^2(\ln\lambda)$ of equation~\eqref{eq:rhoL}
must be elevated to the covariant $g^{\mu\nu}\nabla_\mu\nabla_\nu
(\ln\lambda)$ to be applicable in strong-field and cosmological
regimes.
These targets do not undermine the weak-field predictions presented
here, which are derived in the regime where the approximations hold.

% ── VIII. CONCLUSION ────────────────────────────────────────────────
\section{Conclusion}
\label{sec:conclusion}

We have shown that the cosmological constant problem, the Hubble
tension, and the dark matter problem share a common structural
origin.
When the spacetime metric is sourced by actualization density
$\lambda(x)$ rather than mass-energy density, three consequences
follow without additional parameters: the vacuum does not gravitate
(Theorem~\ref{thm:vacuum}); the Hubble rate is density-dependent,
producing a real physical tension between early- and late-universe
probes; and causally silent particles are categorically prohibited
from gravitating, while baryonic rotation-curve phenomenology is
recovered from the $\rho_\lambda$ term.

The four quantitative predictions of Table~\ref{tab:predictions}
are all testable with current or near-future surveys.
The Eridanus prediction~\eqref{eq:eridanus} is immediately
falsifiable with existing void survey data.
The Einstein radius prediction~\eqref{eq:theta_mw} is testable
with current strong-lensing surveys.
The line-of-sight $H_0$ dependence is testable with DESI and Euclid.

These results flow from the same irreversibility criterion
developed in the companion measurement paper~\cite{Sandeman2026P2}.
The connection to the quantum foundations is direct: off-shell
processes produce no actualization events (they satisfy $\Delta S = 0$),
which is why the vacuum does not gravitate.
The same physical principle that dissolves the measurement problem
also dissolves the cosmological constant problem.

% ── ACKNOWLEDGEMENTS ────────────────────────────────────────────────
\section*{Acknowledgements}
The author acknowledges the use of Anthropic's Claude and Google's
Gemini as AI research partners during the preparation of this
manuscript. The author is solely responsible for all scientific
content and conclusions.

% ── REFERENCES ──────────────────────────────────────────────────────
\begin{thebibliography}{99}

\bibitem{Weinberg1989}
S.~Weinberg,
The cosmological constant problem,
Rev.\ Mod.\ Phys.\ \textbf{61}, 1 (1989).

\bibitem{Planck2020}
Planck Collaboration (N.~Aghanim et al.),
Planck 2018 results VI: Cosmological parameters,
Astron.\ Astrophys.\ \textbf{641}, A6 (2020).

\bibitem{DESI2024}
DESI Collaboration,
DESI 2024 VI: Cosmological constraints from the measurements of
baryon acoustic oscillations,
arXiv:2404.03002 (2024).

\bibitem{Riess2022}
A.~G.~Riess et al.,
A comprehensive measurement of the local value of the Hubble constant
with $1\%$ precision from the HST and the SH0ES team,
Astrophys.\ J.\ Lett.\ \textbf{934}, L7 (2022).

\bibitem{Wong2020}
K.~C.~Wong et al.,
H0LiCOW XIII: A $2.4\%$ measurement of $H_0$ from lensed quasars,
Mon.\ Not.\ Roy.\ Astron.\ Soc.\ \textbf{498}, 1420 (2020).

\bibitem{Verde2019}
L.~Verde, T.~Treu, and A.~G.~Riess,
Tensions between the early and the late Universe,
Nature Astron.\ \textbf{3}, 891 (2019).

\bibitem{Rubin1980}
V.~C.~Rubin, W.~K.~Ford Jr., and N.~Thonnard,
Rotational properties of 21 Sc galaxies with a large range of
luminosities and radii,
Astrophys.\ J.\ \textbf{238}, 471 (1980).

\bibitem{LZ2022}
LZ Collaboration (J.~Aalbers et al.),
First dark matter search results from the LUX-ZEPLIN (LZ) experiment,
Phys.\ Rev.\ Lett.\ \textbf{131}, 041002 (2023).

\bibitem{PandaX2023}
PandaX-4T Collaboration (Y.~Meng et al.),
Dark matter search results from the PandaX-4T commissioning run,
Phys.\ Rev.\ Lett.\ \textbf{127}, 261802 (2021).

\bibitem{Kenworthy2019}
W.~D.~Kenworthy, D.~Scolnic, and A.~Riess,
The local perspective on the Hubble tension: Local structure does
not impact $H_0$,
Astrophys.\ J.\ \textbf{875}, 145 (2019).

\bibitem{Naidoo2016}
K.~Naidoo, L.~Benoit-L\'{e}vy, and O.~Lahav,
The large-scale structure of the Eridanus supervoid,
Mon.\ Not.\ Roy.\ Astron.\ Soc.\ \textbf{459}, L71 (2016).

\bibitem{Clowe2006}
D.~Clowe et al.,
A direct empirical proof of the existence of dark matter,
Astrophys.\ J.\ Lett.\ \textbf{648}, L109 (2006).

\bibitem{Milgrom1983}
M.~Milgrom,
A modification of the Newtonian dynamics as a possible alternative
to the hidden mass hypothesis,
Astrophys.\ J.\ \textbf{270}, 365 (1983).

\bibitem{Verlinde2017}
E.~P.~Verlinde,
Emergent gravity and the dark universe,
SciPost Phys.\ \textbf{2}, 016 (2017).

\bibitem{Poulin2019}
V.~Poulin et al.,
Early dark energy can resolve the Hubble tension,
Phys.\ Rev.\ Lett.\ \textbf{122}, 221301 (2019).

\bibitem{Sandeman2026P2}
J.~F.~Sandeman,
Wave function collapse as irreversible entropy production:
An observer-independent criterion with three experimental predictions,
in preparation (2026).



\bibitem{Sandeman2026P4}
J.~F.~Sandeman,
Relational Actualism: Irreversible Events as the Primitive Basis of
Physical Reality,
in preparation (2026).

\end{thebibliography}

\end{document}
